% 译自 arXiv:2306.13589v2 原 tex 第 3573–3950 行（单元 U09）
\subsection{Serre 对偶与 Fourier--Mukai 变换}\label{ch6.4}


本节包含层上同调与导出函子的若干重要应用。我们从 Serre 建立的著名对偶定理开始，参见 \cite[第 III.7 节]{[Har77]}。

\begin{Lem}
	设 $(X,\mathcal{O}_X)$ 是任意域 $\mathbb{K}$ 上维数为 $n$ 的射影概形（或仅为固有概形）。则 $X$ 有唯一的\textbf{对偶层}\index{sheaf!dualizing} $\mathcal{K}_X\in\textup{obj}(\mathsf{Coh}(X;{}_{\mathcal{O}_X}\!\mathsf{Mod}))$，连同态射 $t:H^n(X,\mathcal{K}_X)\rightarrow\mathbb{K}$，使得对一切 $\mathcal{F}\in\textup{obj}(\mathsf{Coh}(X;{}_{\mathcal{O}_X}\!\mathsf{Mod}))$，自然配对
	\[
	\textup{Hom}_{\mathcal{O}_X}(\mathcal{F},\mathcal{K}_X)\times H^n(X,\mathcal{F})\rightarrow H^n(X,\mathcal{K}_X)\xrightarrow{t}\mathbb{K}
	\]
	给出向量空间的同构 $\theta^0:\textup{Hom}_{\mathcal{O}_X}(\mathcal{F},\mathcal{K}_X)\xrightarrow{\sim} H^n(X,\mathcal{F})^*$（右端表示 $H^n(X,\mathcal{F})$ 的对偶向量空间）。
\end{Lem}

\begin{Thm}[\textbf{Serre 对偶定理}]\index{Serre Duality Theorem}\label{Serreduality}
	设 $(X,\mathcal{O}_X)$ 是 $\mathbb{K}$（代数闭）上维数为 $n$ 的光滑射影概形，因而具有对偶层 $\mathcal{K}_X\in\textup{obj}(\mathsf{Coh}(X;{}_{\mathcal{O}_X}\!\mathsf{Mod}))$。则对一切 $\mathcal{F}\in\textup{obj}(\mathsf{Coh}(X;{}_{\mathcal{O}_X}\!\mathsf{Mod}))$ 与 $k\geq 0$，存在自然的函子性同构
	\begin{equation}
		\theta^k:\textup{Ext}_{\mathcal{O}_X}^k(\mathcal{F},\mathcal{K}_X)\xrightarrow{\sim} H^{n-k}(X,\mathcal{F})^*
	\end{equation}
	（其中 $\theta^0$ 与上一个引理中的同构一致）。

	另一方面，对任意有限秩的 $\mathcal{L}\in\textup{obj}(\mathsf{Sh}^\mathsf{lf}_r(X;{}_{\mathcal{O}_X}\!\mathsf{Mod}))$ 与一切 $k\geq 0$，有 $H^k(X,\mathcal{L})\cong H^{n-k}(X,\mathcal{L}^\vee\otimes_{\mathcal{O}_X}\! \mathcal{K}_X)^*$。
\end{Thm}

% [Tutor Note: 此处极易混淆两个方向。上式的 $\theta^k$ 把 Ext 群映到「上同调的对偶」，次数由 $k$ 翻转为 $n-k$；第二式的 $\mathcal{L}$ 换成 $\mathcal{L}^\vee\otimes\mathcal{K}_X$ 后再取对偶。核对时务必记住 $\mathcal{K}_X$ 是 $n$ 次外幂、其所在次数为 $n$，翻转即由此而来。]

关于射影空间 $X=\mathbb{P}^n$（基域 $\mathbb{K}$）这一更具体的例子，读者可参阅 \cite[定理 III.7.1]{[Har77]}，该结果依赖于 \cite[第 III.5 节]{[Har77]} 中对 $\mathbb{P}^n$ 的层上同调的计算。

\begin{Rem}\label{cansheaf}
	设 $X$ 是代数闭域 $\mathbb{K}$ 上维数为 $n$ 的光滑射影簇。则其对偶层实际上同构于\textbf{典范线丛}\index{canonical line bundle}\footnote{按定义，典范线丛是\textit{微分层} $\Omega_X$ 的 $n$ 次外幂。后者在 \cite[第 II.8 节]{[Har77]} 中构造；在光滑复簇上，它正是 $X$ 上的微分形式层，与我们的直觉相符。}（或\textbf{典范层}\index{sheaf!canonical}）
	\begin{equation}
		K_X\coloneqq \Lambda^n\Omega_X\cong\Omega_X^n\in\text{obj}(\mathsf{Sh}^\mathsf{lf}_1(X;{}_{\mathcal{O}_X}\mathsf{Mod}))\,,
	\end{equation}
	从而是一个可逆凝聚层。

	事实上，记 $\Omega_X^p$ 为 $X$ 上的微分 $p$-形式层，则定理 \ref{Serreduality} 的对偶同构蕴涵
	\begin{equation}\label{Serreisoonvariety}
		H^q(X,\Omega_X^p)\cong H^{n-q}(X,\Omega_X^{n-p})^*\,,
	\end{equation}
	其中 $0\leq p,q\leq n$（只需用到 $\Omega_X^{n-p}\cong (\Omega^p)^\vee\otimes_{\mathcal{O}_X}\! K_X$）。特别地，我们得到 $H^n(X,K_X)\cong\mathbb{K}$，而数 $h^{p,q}(X)\coloneqq\text{dim}(H^q(X,\Omega_X^p))$ 是簇 $X$ 的不变量。

	至此，读者或许已从 \cite[第 7.2 节]{[Imp21]} 中认出了 \eqref{Serreisoonvariety}：若 $X$ 是复流形（即 $\mathbb{C}$ 上的光滑射影簇），这些恒等式可由 Dolbeault 定理 $H_\text{Db}^{p,q}(X)\cong H^q(X,\Omega_X^{p,0})$ 与 Hodge $\star$-算子诱导的同构 $H_\text{Db}^{p,q}(X)\cong H_\text{Db}^{n-p,n-q}(X)$ 联合得到；而诸 $h^{p,q}(X)$ 正是 $X$ 的 Hodge 数。
\end{Rem}

若把注意力集中于簇，Serre 对偶也可以用 \textit{Serre 函子}重述：

\begin{Def}\label{Serrefunc}
	设 $\mathsf{C}$ 是 $\mathbb{K}$-线性范畴。\textbf{Serre 函子}\index{functor!Serre}是一个 $\mathbb{K}$-线性等价 $\mathsf{S}:\mathsf{C}\rightarrow\mathsf{C}$，使得对任意 $X,Y\in\text{obj}(\mathsf{C})$，存在 $\mathbb{K}$-向量空间的函子性同构 $\eta_{X,Y}:\text{Hom}_\mathsf{C}(X,Y)\xrightarrow{\sim}\text{Hom}_\mathsf{C}(Y,\mathsf{S}(X))^*$。

	注意：若 $\mathsf{C}$ 同时是三角范畴，则 $\mathsf{S}$ 是正合函子。
\end{Def}

\begin{Cor}\label{SerredualityviaSerrefunc}
	设 $X=(X,\mathcal{O}_X)\in\textup{obj}(\mathsf{Sch}(\mathbb{K}))$ 光滑、维数为 $n$ 且射影。则（正合）函子
	\begin{equation}
		\mathsf{S}_X\coloneqq(\square\!\otimes_{\mathcal{O}_X}\!K_X)[n]:\mathsf{D^\#}\!(X)\rightarrow\mathsf{D^\#}\!(X)\,,
	\end{equation}
	（对 $\#=\mathsf{b},+,-$，其中 $[n]$ 表示 $\mathsf{D^\#}\!(X)$ 中的 $n$ 次平移）是 $\mathsf{D^\#}\!(X)$ 上的 Serre 函子。
\end{Cor}

\begin{proof}
	我们要证明：对任意两个 $\mathcal{F}^\bullet,\mathcal{G}^\bullet\in\text{obj}(\mathsf{D^b}(X))$，存在定义 \ref{Serrefunc} 中那样的 $\mathbb{K}$-向量空间同构 $\eta_{\mathcal{F}^\bullet\!,\mathcal{G}^\bullet}$。然而这只是定理 \ref{Serreduality} 的直接推论：
	\begin{small}
		\begingroup
		\allowdisplaybreaks\begin{align*}
			\text{Hom}_{\mathsf{D^b}(X)}(\mathcal{G}^\bullet,\,&\mathsf{S}_X(\mathcal{F}^\bullet))^*= & \mkern-20mu{}\\ &=\text{Hom}_{\mathsf{D^b}(X)}(\mathcal{G}^\bullet,(\mathcal{F}^\bullet\otimes_{\mathcal{O}_X}\!K_X)[n])^* & (\text{按定义}) \\
			&\cong\mathsf{R}^n\text{Hom}_{\mathcal{O}_X}^\bullet(\mathcal{G}^\bullet,\mathcal{F}^\bullet\otimes_{\mathcal{O}_X}\!K_X)^* & (\text{由 \eqref{ExtRHomforcomplexes}}) \\
			&\cong\mathsf{R}^n\text{Hom}_{\mathcal{O}_X}^\bullet(\mathsf{R}\mathcal{H}om_{\mathcal{O}_X}(\mathcal{F}^\bullet,\mathcal{G}^\bullet),K_X)^* & (\text{由命题 \ref{lotofequalities}d)}) \\
			&=\text{Ext}_{\mathcal{O}_X}^n(\mathsf{R}\mathcal{H}om_{\mathcal{O}_X}(\mathcal{F}^\bullet,\mathcal{G}^\bullet),K_X)^* & (\text{按定义}) \\
			&\cong H^{n-n}(X,\mathsf{R}\mathcal{H}om_{\mathcal{O}_X}(\mathcal{F}^\bullet,\mathcal{G}^\bullet))^{**} & (\text{由定理 \ref{Serreduality}}) \\
			&\cong \Gamma(X,\mathsf{R}\mathcal{H}om_{\mathcal{O}_X}(\mathcal{F}^\bullet,\mathcal{G}^\bullet))^{**} & (\text{由定义 \ref{sheafcohomology}}) \\
			&\cong \mathsf{R}^0\text{Hom}_{\mathcal{O}_X}^\bullet(\mathcal{F}^\bullet,\mathcal{G}^\bullet)^{**} & (\text{由 \eqref{locHomcompo}}) \\
			&\cong \text{Hom}_{\mathsf{D^b}(X)}(\mathcal{F}^\bullet,\mathcal{G}^\bullet)^{**} & (\text{由 \eqref{ExtRHomforcomplexes}}) \\
			&\cong\text{Hom}_{\mathsf{D^b}(X)}(\mathcal{F}^\bullet,\mathcal{G}^\bullet) & (\text{由 Hom-空间有限性})
		\end{align*}\endgroup\end{small}
	\!\!\!这里我们把 $K_X$ 视为集中在一处的复形（次数为 0 的复形），而援引命题 \ref{lotofequalities} 时，我们实际上用了该命题所列恒等式 $d)$ 的 Hom-对偶版本，并对它再取 $H^n$-上同调。
\end{proof}

% [Tutor Note: 这一串等式的易错处有三。① 第一步把 $\mathsf{S}_X$ 展开为 $(\square\otimes K_X)[n]$，随后用 $\text{Hom}(\mathcal{G}^\bullet,X[n])^*\cong\mathsf{R}^n\text{Hom}^\bullet(\mathcal{G}^\bullet,X)^*$——注意这里 $[n]$ 被「吸收」进了 $\mathsf{R}^n$ 的上标，方向不能反；② 第三行用 Hom-对偶把 $(\mathcal{G}^\bullet,\mathcal{F}^\bullet\otimes K_X)$ 换成 $(\mathsf{R}\mathcal{H}om(\mathcal{F}^\bullet,\mathcal{G}^\bullet),K_X)$ 时，作者「心算」了 $\mathsf{R}\mathcal{H}om$ 的内部 Hom 与张量积的伴随关系，并未写出交换图表；③ 倒数第二步要求 Hom-空间为有限维，这一步在光滑射影情形下成立（相干层上同调的有限性定理），但作者只用括注一笔带过。]

作些微调之后我们看到：要求 $\mathsf{S}_X$ 成为 Serre 函子，必然蕴涵着对一切 $k\geq 0$ 与凝聚层 $\mathcal{F}$ 有 $\textup{Ext}_{\mathcal{O}_X}^k(\mathcal{F},K_X)\cong H^{n-k}(X,\mathcal{F})^*$，而这正是 Serre 对偶的陈述。因此，在光滑射影的 $X\in\textup{obj}(\mathsf{Sch}(\mathbb{K}))$ 上，推论 \ref{SerredualityviaSerrefunc}（限制到凝聚层而非凝聚层复形）与定理 \ref{Serreduality} 等价。

可以借助 Serre 对偶更方便地证明下述生成性结果，它为光滑射影簇的导出范畴给出了一组「正交规范基」：

\begin{Pro}\label{spanningclass}
	设 $X$ 是 $\mathbb{K}$（代数闭）上的光滑射影簇。则\footnote{这里 $k(x)$ 表示闭单点空间 $\{x\}\subset X$ 的常值结构层 $\underline{\mathbb{K}}\in\text{obj}(\mathsf{D^b}(X))$。} $\{k(x)\in\textup{obj}(\mathsf{D^b}(X))\mid x\in X\}$ 是三角范畴 $\mathsf{D^b}(X)$ 的一组\textup{正交规范生成类}，意即下列条件成立：
	\begin{itemize}[leftmargin=0.5cm]
		\item（正交规范性）$\textup{Hom}_{\mathsf{D^b}(X)}(k(x),k(y)[n])\cong\{0\}$，$\forall x\neq y\in X$，$\forall n\in\mathbb{Z}$，且 $\textup{Hom}_{\mathsf{D^b}(X)}(k(x),k(x))\cong\mathbb{K}$，$\forall x\in X$；

		\item（生成性）$\textup{Hom}_{\mathsf{D^b}(X)}(k(x),\mathcal{F}^\bullet[n])\cong\{0\}$，$\forall x\in X$，$\forall n\in\mathbb{Z}$ $\implies$ $\mathcal{F}^\bullet\cong 0^\bullet$。
	\end{itemize}
	（见 \textup{\cite[命题 3.17]{[Huy06]}}。）
\end{Pro}

我们要提醒：这\textit{不}是 \cite[定义 3.22]{[Imp21]} 中经典版本的那个概念：此处我们并非说 $\{k(x)\in\textup{obj}(\mathsf{D^b}(X))\mid x\in X\}$ 构成包含一切有界层复形及其平移与锥的最小全子范畴，从而生成整个 $\mathsf{D^b}(X)$！

% [Tutor Note: 作者的提醒很重要，也是初读最易混的一处。「生成类（spanning class）」是 Huybrechts 书中的判别工具，用于检验函子是否为全忠实（见下文命题 \ref{whenFourierMukaiisequivalence}）；而「生成（generation）」指的是三角范畴意义下由某个类闭合出整个范畴。两者名近而义远，作者在 \cite[定义 3.22]{[Imp21]} 中用的是后者。]

最后，回到概形的范围，我们提及 \cite[第 3.4 节]{[Huy06]} 中讨论的 Serre 对偶之推广。

\begin{Thm}[\textbf{Grothendieck--Verdier 对偶}]\index{Grothendieck--Verdier Duality}\label{GVduality}
	设 $X,Y\in\textup{obj}(\mathsf{Sch}(\mathbb{K}))$ 光滑且射影（或仅为固有），并设 $f:X\rightarrow Y$ 是概形态射，其\textup{相对维数}为 $\textup{dim}(f)\coloneqq\textup{dim}(X)-\textup{dim}(Y)$，\textup{相对对偶丛}为 $K_f\coloneqq K_X\otimes_{\mathcal{O}_X}\!f^*K_Y^\vee\in\textup{obj}(\mathsf{Coh}(X;{}_{\mathcal{O}_X}\!\mathsf{Mod}))$（集中在次数 0）。则对任意 $\mathcal{F}^\bullet\in\textup{obj}(\mathsf{D^b}(X))$ 与 $\mathcal{G}^\bullet\in\textup{obj}(\mathsf{D^b}(Y))$，存在函子性同构
	\begin{equation}
		\mathsf{R}f_*\big(\mathsf{R}\mathcal{H}om_{\mathcal{O}_X}\!\big(\mathcal{F}^\bullet,\mathsf{L}f^*(\mathcal{G}^\bullet)\otimes_{\mathcal{O}_X}\!K_f[\textup{dim}(f)]\big)\big)\cong\mathsf{R}\mathcal{H}om_{\mathcal{O}_Y}\!(\mathsf{R}f_*(\mathcal{F}^\bullet),\mathcal{G}^\bullet)\,.
	\end{equation}
	注意：对两边施加 $\Gamma(Y,\square)$，同时利用 \eqref{globalsectcompo} 与 \eqref{locHomcompo}，取 $H^0$-上同调并最终运用 \eqref{ExtRHomforcomplexes}，便得到
	\[
	\textup{Hom}_{\mathsf{D^b}(X)}(\mathcal{F}^\bullet,\mathsf{L}f^*(\mathcal{G}^\bullet)\otimes_{\mathcal{O}_X}\!K_f[\textup{dim}(f)])\cong \textup{Hom}_{\mathsf{D^b}(Y)}(\mathsf{R}f_*(\mathcal{F}^\bullet),\mathcal{G}^\bullet)\,.
	\]
\end{Thm}

取 $Y=\text{Spec}(\mathbb{K})$ 且 $f$ 为 $X$ 的结构态射（于是 $K_f=K_X$、$\text{dim}(f)=\text{dim}(X)$），我们便得到推论 \ref{SerredualityviaSerrefunc} 的另一证明；若改在凝聚层层面作同样取值，则得到定理 \ref{Serreduality}。

% [Tutor Note: Grothendieck--Verdier 对偶是 Serre 对偶的相对化：当 $Y$ 为单点时，$K_f$ 退化为 $K_X$，四次平移 $[\text{dim}(f)]$ 退化为 $[n]$。核对符号时请注意 $K_f$ 中用的是 $f^*K_Y^\vee$（上标为对偶），且 $K_f$ 被要求「集中在次数 0」——这是作者为避免复形对偶丛的额外复杂度而作的技术性简化，一般情形下 $K_f$ 可以是复形（此时称为相对对偶复形）。]

\vspace*{0.5cm}
\noindent 我们转向另一件联系凝聚层有界导出范畴的重要工具。主要参考文献仍是 \cite{[Huy06]}。

\begin{Def}
	考虑两个光滑射影对象 $X,Y\in\text{obj}(\mathsf{Var}(\mathbb{K}))$ 的积簇 $X\times Y$，带投影 $p_X:X\times Y\rightarrow X$ 与 $p_Y:X\times Y\rightarrow Y$。任取 $\mathcal{P}^\bullet\in\text{obj}(\mathsf{D^b}(X\times Y))$。则与之相伴的\textbf{Fourier--Mukai 变换}\index{Fourier--Mukai transform}是函子\footnote{更形式地，我们应写 $\Phi_{\mathcal{P}^\bullet}(\square)=\mathsf{R}p_{Y*}(\mathsf{L}p_X^*(\square)\overset{\mathsf{L}}{\otimes}\mathcal{P}^\bullet)$；但 $p_X$ 是平坦的，故我们「不必」对它求导（关于平坦态射更多内容见 \cite[第 III.9 节]{[Har77]}）。不过，除非 $\mathcal{P}^\bullet$ 是局部自由层构成的复形，否则我们必须对 $\otimes\equiv\otimes_{\mathcal{O}_{X\times Y}}$ 取左导出。}
	\begin{equation}
		\Phi_{\mathcal{P}^\bullet}\coloneqq\mathsf{R}p_{Y*}(p_X^*(\square)\overset{\mathsf{L}}{\otimes}\mathcal{P}^\bullet):\mathsf{D^b}(X)\rightarrow\mathsf{D^b}(Y)\,,
	\end{equation}
	它因而是正合的，因为正合函子的复合仍正合。我们称 $\mathcal{P}^\bullet$ 为 $\Phi_{\mathcal{P}^\bullet}$ 的\textbf{Fourier--Mukai 核}\index{Fourier--Mukai kernel}；若后者是等价，则称 $X,Y$ 为\textbf{Fourier--Mukai 伙伴}\index{Fourier--Mukai partners}。
\end{Def}

由于 $X\times Y\cong Y\times X$，核 $\mathcal{P}^\bullet$ 同时也定义一个正合函子 $\mathsf{D^b}(Y)\rightarrow\mathsf{D^b}(X)$。我们让两个因子的书写顺序隐含地指示箭头所指的方向。

\begin{Ex}
	Fourier--Mukai 变换其实频繁出没，只是常常不显形。例如，簇的态射 $f:X\rightarrow Y$ 的推前即为 $\mathsf{R}f_*\cong\Phi_{\mathcal{O}_{\Gamma(f)}}:\mathsf{D^b}(X)\rightarrow\mathsf{D^b}(Y)$，即由 $f$ 的图像 $\Gamma(f)\subset X\times Y$ 的结构层 $\mathcal{O}_{\Gamma(f)}\in\text{obj}(\mathsf{D^b}(X\times Y))$（次数 0）所诱导的 Fourier--Mukai 变换。按我们的约定，$\mathcal{O}_{\Gamma(f)}\in\text{obj}(\mathsf{D^b}(Y\times X))$ 反而诱导出拉回 $\mathsf{L}f^*\cong \Phi_{\mathcal{O}_{\Gamma(f)}}:\mathsf{D^b}(Y)\rightarrow\mathsf{D^b}(X)$。更多的实例见 \cite[例 5.4]{[Huy06]} 与 \cite[第 3.4 节]{[Hoc19]}。				\hfill $\blacklozenge$
\end{Ex}

\begin{Rem}
	Fourier--Mukai 变换与 Serre 函子交换。事实上，设 $X,Y$ 为光滑射影簇，$\Phi_{\mathcal{P}^\bullet}:\mathsf{D^b}(X)\rightarrow\mathsf{D^b}(Y)$ 为 $\mathcal{P}^\bullet\in\text{obj}(\mathsf{D^b}(X\times Y))$ 的 Fourier--Mukai 变换。在 $\mathsf{D^b}(X\times Y)$ 中定义有界复形 $\mathcal{P}_\text{L}^\bullet\coloneqq((\mathcal{P}^\bullet)^\vee\otimes p_Y^*K_Y)[\text{dim}(Y)]$ 与 $\mathcal{P}_\text{R}^\bullet\coloneqq((\mathcal{P}^\bullet)^\vee\otimes p_X^*K_X)[\text{dim}(X)]$（它们其实是 $(\mathcal{P}^\bullet)^\vee$ 在投影与 Serre 函子作用下的像）。则作为函子 $\mathsf{D^b}(Y)\rightarrow\mathsf{D^b}(X)$ 有 $\Phi_{\mathcal{P}_\text{L}^\bullet}\cong \Phi_{(\mathcal{P}^\bullet)^\vee}\circ\mathsf{S}_Y$ 与 $\Phi_{\mathcal{P}_\text{R}^\bullet}\cong \mathsf{S}_X\circ\Phi_{(\mathcal{P}^\bullet)^\vee}$，并且二者分别左、右伴随于 $\Phi_{\mathcal{P}^\bullet}$（这本身就意味深长，因为任何正合等价都拥有左伴随与右伴随）。
\end{Rem}

\begin{Lem}
	设 $X,Y,Z$ 为 $\mathbb{K}$ 上的光滑射影簇，$\Phi_{\mathcal{P}_1^\bullet}:\mathsf{D^b}(X)\rightarrow\mathsf{D^b}(Y)$ 与 $\Phi_{\mathcal{P}_2^\bullet}:\mathsf{D^b}(Y)\rightarrow\mathsf{D^b}(Z)$ 分别为 $\mathcal{P}_1^\bullet\in\textup{obj}(\mathsf{D^b}(X\times Y))$ 与 $\mathcal{P}_2^\bullet\in\textup{obj}(\mathsf{D^b}(Y\times Z))$ 的 Fourier--Mukai 变换。
	则 $\Phi_{\mathcal{P}_2^\bullet}\circ \Phi_{\mathcal{P}_1^\bullet}$ 同构于 $\Phi_{\mathcal{P}_3^\bullet}:\mathsf{D^b}(X)\rightarrow\mathsf{D^b}(Z)$，即以下述层为核的 Fourier--Mukai 变换
	\[
	\mathcal{P}_3^\bullet\coloneqq \mathsf{R}p_{X\!Z*}(p_{X\!Y}^*\mathcal{P}_1^\bullet\overset{\mathsf{L}}{\otimes} p_{Y\!Z}^*\mathcal{P}_2^\bullet)\in\textup{obj}(\mathsf{D^b}(X\times Z))\,.
	\]
	（注意：若 $\Phi_{\mathcal{P}_2^\bullet}\circ \Phi_{\mathcal{P}_1^\bullet}$ 不是等价，则 $\mathcal{P}_3^\bullet$ 还可作其他选择。）这使全体核构成一个有良好定义的复合律，其双侧单位元为 $\mathcal{O}_{\Delta_X}\in\textup{obj}(\mathsf{D^b}(X\times X))$，即 $X$ 的对角线的结构层。
\end{Lem}

Orlov 在 \cite{[Orl03]} 中证明了全忠实函子与 Fourier--Mukai 变换之间的下述关系：

\begin{Thm}\label{Orlovthm}
	设 $X,Y$ 为 $\mathbb{K}$ 上的光滑射影簇，$\mathsf{F}:\mathsf{D^b}(X)\rightarrow\mathsf{D^b}(Y)$ 为全忠实正合函子。若 $\mathsf{F}$ 拥有左伴随与右伴随函子，则存在某个 $\mathcal{P}^\bullet\in\textup{obj}(\mathsf{D^b}(X\times Y))$（在同构意义下唯一），使得作为函子有 $\mathsf{F}\cong\Phi_{\mathcal{P}^\bullet}$。此时称 $\mathsf{F}$ \textbf{为 Fourier--Mukai 型}\index{functor!of Fourier--Mukai type}。

	当 $\mathsf{F}$ 是等价时，此条件总成立，且此时底层的簇满足 $\textup{dim}(X)=\textup{dim}(Y)$。特别地，$\textup{Auteq}(\mathsf{D^b}(X))$ 中每个自等价都由某个核 $\mathcal{P}^\bullet\in\textup{obj}(\mathsf{D^b}(X\times X))$ 诱导，从而 $\textup{Auteq}(\mathsf{D^b}(X))$ 具有群的结构。
\end{Thm}

% [Tutor Note: 作者写「此条件总成立」时省略了证明。Orlov 的原始论证（\cite{[Orl03]}，亦见 \cite[定理 5.14]{[Huy06]}）分两步：先由全忠实性构造伴随的「左右双向」结构，再用 Bondal--Orlov 的表示定理把正合函子在 $X$ 的生成类 $\{k(x)\}$ 上的取值组装成一个核。此处关键的技术假设是 $\mathsf{D^b}(X)$ 具生成类且 $X$ 光滑射影，故任何「逐点」给定的数据都能拼成 $\mathsf{D^b}(X\times Y)$ 中的对象。]

另有一些著称的结果，其证明取自 Fourier--Mukai 理论：

\begin{Pro}\label{whenprojvarareequiv}
	设 $X,Y$ 为 $\mathbb{K}$ 上的光滑射影簇。则：
	\begin{itemize}[leftmargin=0.5cm]
		\item 若存在等价 $\mathsf{F}:\mathsf{Coh}(X;{}_{\mathcal{O}_X}\!\mathsf{Mod})\rightarrow \mathsf{Coh}(Y;{}_{\mathcal{O}_Y}\!\mathsf{Mod})$，则作为簇有 $X\cong Y$。

		\item 若存在正合等价 $\mathsf{F}:\mathsf{D^b}(X)\!\rightarrow\! \mathsf{D^b}(Y)$——即 $X, Y$ 导出等价——且 $X$ 的（反）典范线丛是丰富的，则 $X\cong Y$（并且 $Y$ 的（反）典范线丛也是丰富的）。\footnote{除去椭圆曲线，这一由 Bondal 与 Orlov 得到的结果（见 \cite{[BO97]}）完全分类了维数为 1 的光滑射影簇的导出范畴！二人还在同样假定丰裕性的前提下证明了 $\text{Auteq}(\mathsf{D^b}(X))\cong\mathbb{Z}\times(\text{Aut}(X)\ltimes\text{Pic}(X))$——其中 $\mathbb{Z}$ 经平移 $[1]$ 作用、$\text{Aut}(X)$ 经拉回作用、$\text{Pic}(X)$ 经张量积作用——这实际上确定了 $\mathsf{D^b}(X)$ 的复杂程度。}
	\end{itemize}
\end{Pro}

值得一问的是：Fourier--Mukai 变换何时是全忠实函子、甚至等价？Bondal、Orlov 与 Bridgeland（见 \cite{[Bri98]}）给出如下回答。

\begin{Pro}\label{whenFourierMukaiisequivalence}
	设 $X,Y$ 为 $\mathbb{K}$ 上的光滑射影簇，$n=\textup{dim}(X)$，$\Phi_{\mathcal{P}^\bullet}:\mathsf{D^b}(X)\rightarrow\mathsf{D^b}(Y)$ 为核 $\mathcal{P}^\bullet\in\textup{obj}(\mathsf{D^b}(X\times Y))$ 的 Fourier--Mukai 变换。则：
	\begin{itemize}[leftmargin=0.5cm]
		\renewcommand{\labelitemi}{\textendash}
		\item $\Phi_{\mathcal{P}^\bullet}$ 全忠实，当且仅当对任意两个闭点 $x_1,x_2\in X$ 成立
		\[\begin{small}
			\textup{Hom}_{\mathsf{D^b}(Y)}\big(\Phi_{\mathcal{P}^\bullet}(k(x_1)),\Phi_{\mathcal{P}^\bullet}(k(x_2))[l]\big)\!\cong\!
			\begin{cases}
				\mathbb{K} & \text{若 }\; x_1\!=\!x_2 \;\wedge\; l\!=\!0 \\
				\{0\} & \text{若 }\; x_1\!\neq\! x_2 \;\vee\; l\!<\!0\;\vee\; l\!>\!n
			\end{cases}\,;
		\end{small}\]

		\item 若 $\Phi_{\mathcal{P}^\bullet}$ 已经全忠实，则它还是等价，当且仅当 $\textup{dim}(X)=\textup{dim}(Y)$ 且 $\mathcal{P}^\bullet\overset{\mathsf{L}}{\otimes} p_X^*K_X\cong \mathcal{P}^\bullet\overset{\mathsf{L}}{\otimes} p_Y^*K_Y$；等价地，当且仅当对一切闭点 $x\in X$ 有 $\Phi_{\mathcal{P}^\bullet}(k(x))\otimes_{\mathcal{O}_Y}\!K_Y\cong \Phi_{\mathcal{P}^\bullet}(k(x))$。

		\item 若 $\textup{dim}(X)=\textup{dim}(Y)$ 且典范线丛平凡，即 $K_X\cong\mathcal{O}_X$、$K_Y\cong\mathcal{O}_Y$，则任一全忠实正合函子 $\mathsf{D^b}(X)\rightarrow\mathsf{D^b}(Y)$ 自动是等价。
	\end{itemize}
\end{Pro}

% [Tutor Note: 原文此处三个 case 分支写作 `$if $` 开头的伪代码（形如 `$\mathbb{K}$ & $if $\; x_1\!=\!x_2 ...`），那是排版瑕疵而非数学内容；译文正好借机改排为 \text{若\ldots} 以符合中文，公式本身未动。核对符号时请注意：第二条件的 $l$ 范围是 $0\leq l\leq n$，即**不允许负次数**，这与「正交规范」中 $l$ 取遍 $\mathbb{Z}$ 的直觉不同。]

还有其他若干判别 Fourier--Mukai 变换性质的准则。好奇的读者可参阅 \cite[第 7 章]{[Huy06]}。


\subsection{Calabi--Yau 簇的导出范畴}\label{ch6.5}


为定义 Calabi--Yau 流形上凝聚层的有界导出范畴，我们必须先弄清楚如何把后者视为簇、并最终视为概形。由 \cite[第 7.2 节]{[Imp21]} 我们知道：

\begin{Def}\label{Calabiyauman}
	一个 \textbf{Calabi--Yau 流形}\index{Calabi--Yau manifold} $X^n=(X,\varpi)$ 是紧 Kähler 流形 $(X,J,\kappa^\mathbb{C})$，其黎曼和乐群 $\text{Hol}(\kappa^\mathbb{C})=\text{SU}(n)$，并且配备一个（相差一个相位意义下）唯一的全纯体积形式 $\varpi\in\Omega_X^{n,0}(X)\cong\mathcal{O}_X(X)$，称为 \textit{Calabi--Yau 形式}，满足等式 $\kappa^n/n!=(-1)^{n(n-1)/2}(i/2)^n\varpi\wedge\overline{\varpi}$。
\end{Def}

% [Tutor Note: 此处「和乐群 $=\text{SU}(n)$」是 Calabi--Yau 的等价刻画之一（而非仅「$\subseteq\text{SU}(n)$」）：若和乐群真含于 $\text{SU}(n)$，则流形或是超 Kähler（和乐 $\text{Sp}(n/2)$）、或是复环面等特例，仍属广义 CY 但并非本文所定义的严格意义。另注意体积形式的归一化等式 $\kappa^n/n!=(-1)^{n(n-1)/2}(i/2)^n\varpi\wedge\overline{\varpi}$ 中的相位与符号约定依作者而异（常见写法用 $(-1)^{n(n-1)/2}(\tfrac{i}{2})^n\varpi\wedge\overline{\varpi}$，亦有吸收到 $\varpi$ 的相位里者），核对时以原文为准。]

一个既有启发性、又带警示意味的观察是：$\mathbb{C}$ 上的光滑（抽象）簇可以解释为复流形，但反之一般不真。让我们简要看看这是怎么发生的，以及 Calabi--Yau 流形是否有希望逆转这一蕴涵（剧透：它们确实可以）。我们参阅 \cite[附录 B]{[Har77]} 与 \cite{[Vak17]}。

\begin{Def}
	记复单位圆盘为 $\mathbb{D}^n\coloneqq\{z=(z_1,...,z_n)\in\mathbb{C}^n\mid |z_i|<1\;\forall 1\leq i\leq n\}\subset\mathbb{C}^n$，并配以全纯函数的结构层。一个 \textbf{复解析空间}\index{complex analytic space}是局部环空间 $(\mathfrak{X},\mathcal{O}_\mathfrak{X})$，它容许一个开覆盖，使得覆盖中的每个集合都（作为局部环空间）同构于某个\textbf{解析簇}\index{variety!analytic} $Y\coloneqq V(\{f_1,...,f_k\})\subset\mathbb{D}^n$——即有限个全纯函数 $f_1,...,f_k:\mathbb{D}^n\rightarrow\mathbb{C}$ 的零点轨迹，从而是（关于标准拓扑 $\mathcal{T}_{\mathbb{C}^n}$ 的）闭子集，并赋予 $\mathcal{O}_Y\coloneqq\mathcal{O}_{\mathbb{D}^n}/\langle\{f_1,...,f_k\}\rangle$ 而成为局部环空间。\footnote{形式地说，复解析空间即是 $\mathbb{C}$ 上的一个局部环空间，它在局部同构于刚才定义的某个\textit{模型空间} $(Y,\mathcal{O}_Y)$；所说的同构也称为 \textit{$X$ 的局部模型}。} 取复解析空间之间的态射为底层局部环空间的态射，便得到复解析空间范畴 $\mathsf{Sp}_\mathbb{C}^\text{an}$。

	现在设 $(X,\mathcal{O}_X)$ 是 $\mathbb{C}$ 上的有限型概形。则其\textbf{相伴复解析空间}\index{complex analytic space!associated}（或称\textit{复解析簇}）$(X_\text{an},\mathcal{O}_{X_\text{an}})$ 如下得到：按定义，$X$ 容许由形如 $Y_i=\text{Spec}\big(\mathbb{C}[x_1,...,x_n]/\langle\{f_{i,1},...,f_{i,k}\}\rangle\big)$ 的仿射概形构成的开覆盖，其多项式可以看作全纯函数 $\mathbb{C}^n\rightarrow\mathbb{C}$，因而定义了复解析子空间 $(Y_i)_\text{an}\subset\mathbb{C}^n$；于是，按诸 $Y_i$ 粘合成 $X$ 的同样方式把诸 $(Y_i)_\text{an}$ 相容地粘合，便得到 $X_\text{an}$，其结构层 $\mathcal{O}_{X_\text{an}}$ 同理。这一构造是函子性的，于是我们得到所谓的\textbf{解析化函子}\index{functor!analytification}
	\[
	\square_\text{an}:\{X\in\text{obj}(\mathsf{Sch}(\mathbb{C}))\mid X\text{ of finite type}\}\rightarrow \mathsf{Sp}_\mathbb{C}^\text{an}\,.
	\]

	此外，回想：任一凝聚 $\mathcal{O}_X$-模 $\mathcal{F}\in\text{obj}(\mathsf{Coh}(X;{}_{\mathcal{O}_X}\!\mathsf{Mod}))$ 按定义 \ref{coherentonringed} 局部地形如 $\mathcal{F}|_U\cong\mathcal{C}oker(\psi)$（限制到任一 Zariski 开集 $U\subset X$ 上），其中 $\psi$ 是有限生成自由层之间的态射，落在某个正合序列 $\mathcal{O}_X|_U^{\oplus n_i}\xrightarrow{\psi}\mathcal{O}_X|_U^{\oplus m_i}\xrightarrow{\chi}\mathcal{F}|_U\rightarrow 0$ 之中。于是 $\mathcal{F}$ 可辨认出一个\textbf{凝聚解析层}\index{sheaf!coherent analytic} $\mathcal{F}_\text{an}$：把由诱导态射 $\psi_\text{an}:(\mathcal{O}_X)_\text{an}|_{U_\text{an}}^{\oplus n_i}\rightarrow(\mathcal{O}_X)_\text{an}|_{U_\text{an}}^{\oplus m_i}$ 而来的 $\mathcal{F}_\text{an}|_{U_\text{an}}\cong\mathcal{C}oker(\psi_\text{an})$，按各个 $\mathcal{F}|_U$ 粘成 $\mathcal{F}$ 的同样方式粘合而成。凝聚解析 $\mathcal{O}_{X_\text{an}}$-模构成范畴 $\mathsf{Coh}(X_\text{an};{}_{\mathcal{O}_{X_\text{an}}}\!\!\!\mathsf{Mod})$。
\end{Def}

对 $X\in\text{obj}(\mathsf{Sch}(\mathbb{C}))$，已知下述对应：$X$ 分离当且仅当 $X_\text{an}$ 是 Hausdorff 的；$X$ 既约当且仅当 $X_\text{an}$ 既约；$X$ 光滑当且仅当 $X_\text{an}$ 是复流形；$X$ 固有当且仅当 $X_\text{an}$ 紧。

我们要指出：底层拓扑空间之间的隐含映射 $\varphi:X_\text{an}\rightarrow X$ 把 $X_\text{an}$ 双射到 $X$ 的闭点集上，并可提升为局部环空间的态射，从而对任一 $\mathcal{F}\in\text{obj}(\mathsf{Coh}(X;{}_{\mathcal{O}_X}\!\mathsf{Mod}))$ 有 $\mathcal{F}_\text{an}\cong\varphi^*\mathcal{F}$；此外，对每个 $n\in\mathbb{Z}$，层上同调群之间存在自然映射 $H^n(X,\mathcal{F})\rightarrow H^n(X_\text{an},\mathcal{F}_\text{an})$。

% [Tutor Note: 上面这串对应关系中，「$X$ 固有 iff $X_\text{an}$ 紧」正是把代数几何的「固有性」与拓扑的「紧性」沟通的桥梁，也是 GAGA 型定理的核心技术点之一。而「$X$ 光滑 iff $X_\text{an}$ 是复流形」解释了为什么在 $\mathbb{C}$ 上可以把光滑簇直接当作复流形来算 Dolbeault 上同调。注意 $\mathcal{F}_\text{an}\cong\varphi^*\mathcal{F}$ 用的是 $\varphi$ 的**拉回**而非推前，这与 $\varphi$ 的连续像并不构成同胚（$\varphi$ 只是双射到闭点上，拓扑可能不同）。]

如今，关于簇（按定义 \ref{abstractvar} 的意义）与复流形是否可互换的问题，取决于 $\square_\text{an}$ 究竟有多好：由上可知，「光滑抽象簇可以解释为复流形」不过是赋值 $X\mapsto X_\text{an}$；至于反方向，我们需要 $\square_\text{an}$ 作为函子是本质满的。可惜这并不成立……但若取 $X$ 为射影的（这蕴涵 $X$ 在 $\mathbb{C}$ 上有限型），则解析化便使我们的梦想成真，此即 Serre 在 \cite{[Ser56]} 中所证。

\begin{Thm}
	特殊化的解析化函子
	\begin{equation}\label{analytfunc}
		\mkern-12mu\square_\textup{an}\!:\!\{X\!\in\!\textup{obj}(\mathsf{Sch}(\mathbb{C}))\mid X\text{ projective}\}\rightarrow \{\mathfrak{X}\!\in\!\textup{obj}(\mathsf{Sp}_\mathbb{C}^\textup{an})\mid \mathfrak{X}\text{ compact}\}
	\end{equation}
	是本质满的\footnote{具体说来，若 $\mathfrak{X}$ 是复流形 $\mathbb{C}P^n$ 的紧解析子空间（这总是可以做到的），则存在子概形 $X\subset\mathbb{C}P^n$ 使得 $X_\text{an}\cong\mathfrak{X}$。}，且在对象上是同构单射（后者意指：$X_\textup{an}\!\cong\! X_\textup{an}'\!\implies\! X\!\cong\! X'$）。

	此外，其在凝聚层上诱导的函子
	\begin{equation}
		\square_\textup{an}:\mathsf{Coh}(X;{}_{\mathcal{O}_X}\!\mathsf{Mod})\rightarrow \mathsf{Coh}(X_\textup{an};{}_{\mathcal{O}_{X_\textup{an}}}\!\!\!\mathsf{Mod})
	\end{equation}
	是等价，从而是本质满的，且在对象上同构单射，这就使自然映射
	\[
	H^n(X,\mathcal{F})\rightarrow H^n(X_\textup{an},\mathcal{F}_\textup{an})
	\]
	对一切 $n\in\mathbb{Z}$ 成为同构。
\end{Thm}

% [Tutor Note: 这就是 Serre 的 GAGA 定理（\cite{[Ser56]}，\textit{Géométrie algébrique et géométrie analytique}）。作者只写「如 Serre 所证」，省略了全部证明：其核心是把射影空间上的凝聚层用扭层 $\mathcal{O}(k)$ 的有限消解表出，再对复解析版本重复同一消解，从而把上同调的计算从「解析世界」拉回「代数世界」。应注意定理要求**射影**：去掉射影性后本质满性立刻失败（存在非代数化的紧复流形）。]

那么，若把 \eqref{analytfunc} 的靶空间扩大为紧复流形 $\mathfrak{X}$，本质满性是否仍然成立？若如此，即存在某射影概形 $X$ 使 $X_\text{an}\cong\mathfrak{X}$（事实上唯一），则称 $\mathfrak{X}$ \textbf{射影代数的}\index{complex analytic space!projective algebraic}。由 \cite[第 B.3 节]{[Har77]} 我们学到：

\begin{Thm}
	\textup{（Riemann）} 每个 1 维紧复流形，换言之每个紧 Riemann 曲面——包括 Calabi--Yau 1-维流形——都是射影代数的。

	\textup{（Chow--Kodaira）} 每个 2 维、且带有两个代数无关亚纯函数的紧复流形——例如 Calabi--Yau 2-维流形——都是射影代数的。\footnote{事实上，此时它是 \textit{Moishezon 流形}，而按 \cite[命题 B.3.3]{[Har77]}，这是必要条件。}
\end{Thm}

在维数大于 2 时，对一般的紧复流形而言情形变得更为复杂。话虽如此，据 \cite[定理 5.8]{[GHJ03]}，我们运气不错：

\begin{Thm}\label{3+DCYareprojectivealgebraic}
	任一维数 $n\geq 3$ 的 Calabi--Yau 流形，都作为复流形同构于复射影空间 $\mathbb{C}P^m$（$m>n$）的某子流形，而该子流形是 $\mathbb{C}$ 上某个射影概形的解析化，从而此流形是射影代数的。\footnote{更一般地，任何 Moishezon 的 Kähler 流形都是射影代数的。}\footnote{另一方面，存在大量并非代数的 $K3$ 曲面——它们是 Calabi--Yau 2-维流形的实例——因而不能嵌入复射影空间。还应注意，维数 $\geq 2$ 的复环面也可能出现同样的问题（这提示对定理 \ref{3+DCYareprojectivealgebraic} 值得进一步探究，本文未予深论；感谢 Biran 教授提出这一点）。}
\end{Thm}

因此我们看到：当 $\square_\text{an}$ 取值于「本身是 Calabi--Yau 流形的紧复解析空间」这一范畴时，它是本质满的。映射到这些解析空间的射影概形，恰是所谓的 \textit{Calabi--Yau 簇}，这一概念也可以在任意代数闭域 $\mathbb{K}$ 的更广语境中定义：

\begin{Def}\label{Calabiyauvar}
	代数闭域 $\mathbb{K}$ 上的 $n$ 维 \textbf{Calabi--Yau 簇}\index{Calabi--Yau variety}，是指光滑固有簇 $X=(X,\mathcal{O}_X)\in\text{obj}(\mathsf{Sch}(\mathbb{K}))$，其典范线丛 $K_X$（见注 \ref{cansheaf}）平凡，即 $K_X\cong\mathcal{O}_X$（从而是 $X$ 上正则函数层），且对每个 $0< j<n$ 其层上同调群平凡：$H^j(X,\mathcal{O}_X)\cong\{0\}$。

	在 $\mathbb{K}=\mathbb{C}$ 时，我们只需把「固有」簇换成「射影」簇（由命题 \ref{whenproper}，射影蕴涵固有），并把「正则」函数换成「全纯」函数（见例 \ref{sheafholofcts}）。
\end{Def}

\begin{Rem}
	与 \cite[注 7.7]{[Imp21]} 中所讨论的流形情形类似，数学文献对 Calabi--Yau 簇常采用彼此可能不等价的定义：例如，有人放弃关于层上同调群平凡的要求（于是，当该要求也满足时，人们谈论的是\textit{真} Calabi--Yau 簇）；而许多作者只要求第一 Chern 类\footnote{经由 Čech 上同调可以证明 $\text{Pic}(X)\cong H^1(X,\mathcal{O}_X^\times)$，其中 $\text{Pic}(X)$ 如定义 \ref{Picardgroup} 所示，$\mathcal{O}_X^\times$ 是 $X$ 上\textit{非零}正则函数层。此外，阿贝尔层短正合序列 $0\rightarrow\mathbb{Z}\rightarrow\mathcal{O}_X\xrightarrow{\text{exp}}\mathcal{O}_X^\times\rightarrow 0$ 诱导出带连接同态的长正合序列，而当 $X$ 射影时，该连接同态恰是 $c_1:H^1(X,\mathcal{O}_X^\times)\cong\text{Pic}(X)\rightarrow H^2(X,\mathbb{Z})$。这给出了第一 Chern 类 $c_1$ 的一个层论定义。此论证在 \cite[第 4.5.1 节]{[ABC+09]} 中被强调。} $c_1(X)$ 为零，而这又蕴涵典范线丛平凡。

	人们甚至可以更一般地在非代数闭、或特征非零的域上谈论 Calabi--Yau 簇，但此处我们无意沿此路径深入。
\end{Rem}

% [Tutor Note: 这里存在一个真实的文献分歧，切勿死记一种定义。本文采用的是「$K_X$ 平凡 + 中间上同调全灭」的**强**定义（常称 strict / true CY）；多数文献（含 Huybrechts 的 \cite{[Huy06]}）只要求 $K_X$ 平凡。二者在维数 1、2 时差别不大，维数 $\geq 3$ 时则确有反例。作者在下文 \ref{dercatofCY} 之后明确说「我们偏爱的定义」——因其载有最多信息。】

拆解第 \ref{ch5.5} 节讨论过的概形的诸多属性之后，定义 \ref{Calabiyauvar} 表明：Calabi--Yau 簇 $X\in\text{obj}(\mathsf{Sch}(\mathbb{K}))$ 是光滑的\textit{完备}簇，且总是在定义 \ref{abstractvar} 的抽象意义下如此；因此由图表 \eqref{bigdiagram}，它表明：
\vspace*{0.3cm}

\minibox[frame]{Calabi--Yau 簇是 $\mathbb{K}$ 上的光滑、整、诺特、分离的有限型概形，\\并附加「典范线丛平凡」以及「大部分层上同调群平凡」两项要求。}\vspace*{0.3cm}

当然，$X$ 经由函子 \eqref{Xi} 辨认出某个真正的簇 $X\in\text{obj}(\mathsf{Var}(\mathbb{K}))$，\footnote{此时由定义 \ref{smoothsch}，$X$ 作为簇也光滑，因为它是整的，从而是既约且不可约的。}它同时也是一个配以正则函数层的环空间。

我们再次强调：各家作者对 Calabi--Yau 簇的代数定义常不一致。我们所偏爱、且第二部分的理论据以构建的那个定义，有一个特别的优点：它把我们所可能企望的信息量最大化了——这一点稍后就会看到。

万事俱备，现在让我们最后一次描述 $\mathsf{D^b}(X)$ 的结构，其方式一如定义 \ref{dercatdef}，只是这次假定 $X$ 是 Calabi--Yau 流形，并因而隐含地假定了相应的复 Calabi--Yau 簇 $X=(X,\mathcal{O}_X)\in\text{obj}(\mathsf{Sch}(\mathbb{C}))$。

\begin{Def}\label{dercatofCY}
	设 $X$ 是 Calabi--Yau 流形，其相伴阿贝尔范畴为 $\mathsf{Coh}(X;{}_{\mathcal{O}_X}\!\mathsf{Mod})$，该范畴的同伦范畴\footnote{这里我们把定义 \ref{dercatofsch} 中采用的记号推广到 $\mathsf{Coh}(X;{}_{\mathcal{O}_X}\!\mathsf{Mod})$ 的（有界）同伦范畴以及（有界）复形范畴。} $\mathsf{H^b}(X)\!\coloneqq\!\mathsf{H^b}(\mathsf{Coh}(X;{}_{\mathcal{O}_X}\!\mathsf{Mod}))$ 的局部化类为 $S\coloneqq\{\text{quasi-isomorphisms in }\mathsf{H^b}(X)\}$。则 $X$ 的有界导出范畴 $\mathsf{D^b}(X)=\mathsf{D^b}(\mathsf{Coh}(X;{}_{\mathcal{O}_X}\!\mathsf{Mod}))=\mathsf{H^b}(X)[S^{-1}]$ 具有如下结构：
	\begin{itemize}[leftmargin=0.5cm]
		\renewcommand{\labelitemi}{\textendash}
		\item $\text{obj}(\mathsf{D^b}(X))$ 中的对象是有界上链复形 $(\mathcal{F}^\bullet,d_\mathcal{F}^\bullet)$，其中每个 $\mathcal{F}^n\in\text{obj}(\mathsf{Coh}(X;{}_{\mathcal{O}_X}\!\mathsf{Mod}))$ 是凝聚 $\mathcal{O}_X$-模；

		\item 态射 $\varphi\in\text{Hom}_{\mathsf{D^b}(X)}(\mathcal{F}^\bullet,\mathcal{G}^\bullet)$ 是「屋顶」的等价类 $[(\psi^\bullet,\chi^\bullet)]$
		\begin{equation}
			\begin{tikzcd}
				& \mathcal{F}'^\bullet\arrow[dl, "\psi^\bullet"']\arrow[dr, "\chi^\bullet"] & \\
				\mathcal{F}^\bullet\arrow[rr, squiggly, "\varphi"'] & & \mathcal{G}^\bullet
			\end{tikzcd}\quad,
		\end{equation}
		其中 $\chi^\bullet\in\text{Hom}_{\mathsf{H^b}(X)}(\mathcal{F}'^\bullet,\mathcal{G}^\bullet)$ 是凝聚层复形的态射，且 $\psi^\bullet\in S$。两个屋顶等价，是指它们能被置于 \eqref{equivroofs} 那样的公共屋顶之下；而态射的复合按 \eqref{roofcompo} 进行。
	\end{itemize}
	导出范畴配有一个局部化函子 $\mathsf{L}:\mathsf{H^b}(X)\rightarrow\mathsf{D^b}(X)$，它在对象上是恒等，并把任一 $\chi^\bullet\in\text{Hom}_{\mathsf{H^b}(X)}(\mathcal{F}^\bullet,\mathcal{G}^\bullet)$ 映为类 $[(\text{id}_\mathcal{F}^\bullet,\chi^\bullet)]\in\text{Hom}_{\mathsf{D^b}(X)}(\mathcal{F}^\bullet,\mathcal{G}^\bullet)$。此外，对任一 $\psi^\bullet\in S$，$\mathsf{L}(\psi^\bullet)$ 是同构；且若另有函子 $\mathsf{F}:\mathsf{H^b}(X)\rightarrow\mathsf{D}$ 也如此作用，则它可分解为 $\mathsf{F}=\mathsf{G}\circ\mathsf{L}$，其中 $\mathsf{G}:\mathsf{D^b}(X)\rightarrow\mathsf{D}$ 唯一。
\end{Def}

借此机会把本次综述中遇到的最重要结果汇集起来，以下是我们所能推知的 $\mathsf{D^b}(X)$ 基本特征的总结：

\begin{itemize}[leftmargin=0.5cm]
\item 由推论 \ref{D(A)istriangulated}，当配以定义 \ref{translfunc} 的平移函子 $T=T[1]:\mathsf{D^b}(X)\rightarrow \mathsf{D^b}(X)$ 以及定义 \ref{exacttriangles} 中复形的区分三角形时，$\mathsf{D^b}(X)$ 是三角范畴（见定义 \ref{triangcat}）。
\newline 此外，对任一 $\mathcal{F}^\bullet\in\text{obj}(\mathsf{D^b}(X))$，函子 $\text{Hom}_{\mathsf{D^b}(X)}(\mathcal{F}^\bullet,\square)$、$\text{Hom}_{\mathsf{D^b}(X)}(\square,\mathcal{F}^\bullet):\mathsf{D^b}(X)\rightarrow\mathsf{Ab}$ 都是上同调的／三角的（由命题 \ref{Homareexact}），上同调函子 $H^n:\mathsf{D^b}(X)\rightarrow\mathsf{Coh}(X;{}_{\mathcal{O}_X}\!\mathsf{Mod})$ 亦然（由定理 \ref{longexactcohomtriangles}）；因此 $\mathsf{D^b}(X)$ 中的区分三角形诱导出长的正合 Hom-序列／上同调序列。

\item 命题 \ref{AinD(A)} 中的全忠实函子 $\mathsf{J}':\mathsf{Coh}(X;{}_{\mathcal{O}_X}\!\mathsf{Mod})\hookrightarrow\mathsf{D^b}(X)$ 使我们能把 $\mathsf{Coh}(X;{}_{\mathcal{O}_X}\!\mathsf{Mod})$ 解释为 $\mathsf{D^b}(X)$ 中一切 $H^0$-复形构成的全子范畴：凝聚 $\mathcal{O}_X$-模 $\mathcal{F}$ 与其 0-复形 $\mathcal{F}[0]^\bullet$ 等同，且对每对 $\mathcal{O}_X$-模 $\mathcal{F}$、$\mathcal{G}$ 有 $\text{Hom}_{\mathcal{O}_X}(\mathcal{F},\mathcal{G})\cong\text{Hom}_{\mathsf{D^b}(X)}(\mathcal{F}[0]^\bullet,\mathcal{G}[0]^\bullet)$。
\newline 请注意：$\mathsf{D^b}(X)$ 本身\textit{不是}阿贝尔范畴，但由引理 \ref{D(A)additive} 它是加性范畴（从而例如可借助定理 \ref{YonedaEmbThm} 来研究）。

\item 由定理 \ref{Extprop} 与注 \ref{HomExtRem}，\eqref{Extmodules} 中定义的凝聚 $\text{Ext}$-模把任一凝聚 $\mathcal{O}_X$-模的短正合序列映为以 $\text{Hom}_{\mathcal{O}_X}$-群开始、并在次数 $n>\text{dim}(X)$ 上平凡的有界正合序列，因为当 $X$ 射影时 $\mathsf{Coh}(X;{}_{\mathcal{O}_X}\!\mathsf{Mod})$ 的同调维数 \eqref{homologicaldimension} 为 $\text{dim}(X)$。

\item 对 Calabi--Yau 的 $X$，阿贝尔范畴 $\mathsf{QCoh}(X;{}_{\mathcal{O}_X}\!\mathsf{Mod})$ 有足够多的内射对象（见注 \ref{QCohhasenoughinjectives}），从而由命题 \ref{globalsectforcoherentbdd}，$\Gamma(X,\square):\mathsf{QCoh}(X;{}_{\mathcal{O}_X}\!\mathsf{Mod})\rightarrow \mathsf{Vect}_\mathbb{K}$ 的右导出函子存在。
\newline 特别地，层上同调在 $X$ 上有良好定义，并可用 Čech 上同调计算（由定理 \ref{sheafcohomthroughCech}，该定理适用是因为 $X$ 诺特且分离）。给定 Calabi--Yau 簇的仿射态射 $f:X\rightarrow Y$，对一切凝聚 $\mathcal{O}_X$-模 $\mathcal{F}$ 与 $n\geq 0$ 有 $H^n(X,\mathcal{F})\cong H^n(Y,f_*\mathcal{F})$（此即推论 \ref{sheafcohomthroughCechRem}）。

\item 仍然借助 $\mathsf{Qcoh}(X;{}_{\mathcal{O}_X}\!\mathsf{Mod})$ 与引理 \ref{boundedCohinQCoh}，我们可以完成第 \ref{ch6.3} 节中关于 $\mathsf{D^b}(X)$ 上正合导出函子之存在性的全部证明（因为 $X$ 满足我们在那里所作的一切假设）。设 $f:X\rightarrow Y$ 是 Calabi--Yau 簇的态射，则我们有 $\mathsf{R}f_*:\mathsf{D^b}(X)\rightarrow\mathsf{D^b}(Y)$（当 $f$ 射影或固有）与 $\mathsf{L}f^*:\mathsf{D^b}(Y)\rightarrow\mathsf{D^b}(X)$。我们还可以考虑 $\mathsf{R}\mathcal{H}om_{\mathcal{O}_X}^\bullet(\square,\square):\mathsf{D^b}(X)^\text{opp}\times\mathsf{D^b}(X)\rightarrow\mathsf{D^b}(X)$，连同 $\mathsf{R}\text{Hom}_{\mathcal{O}_X}^\bullet(\square,\square):$ $\mathsf{D^b}(X)^\text{opp}\times\mathsf{D^b}(X)\rightarrow\mathsf{D^b}(\mathsf{Vect}_\mathbb{K})$（由此得到局部的 $\mathcal{E}xt$-层、对偶有界复形与 $\text{Ext}$-模），以及 $\square\overset{\mathsf{L}}{\otimes}_{\mathcal{O}_X}\!\square:\mathsf{D^b}(X)\times\mathsf{D^b}(X)\rightarrow\mathsf{D^b}(X)$。所有这些都满足命题 \ref{lotofequalities} 所描述的恒等式。

\item 由定理 \ref{Serreduality}，Serre 对偶在 Calabi--Yau 的 $X$ 上同样成立，于是对一切 $k\geq 0$ 得到同构 $\text{Ext}_{\mathcal{O}_X}^k(\mathcal{F},\mathcal{O}_X)\cong H^{n-k}(X,\mathcal{F})^*$；而推广之的 Grothendieck--Verdier 对偶（定理 \ref{GVduality}）对 Calabi--Yau 簇之间的任一态射 $f:X\rightarrow Y$ 也成立。此外，命题 \ref{spanningclass} 为 $\mathsf{D^b}(X)$ 提供了一组正交规范生成类。
\newline Fourier--Mukai 理论适用于 Calabi--Yau 簇：定理 \ref{Orlovthm} 说明，两个给定 Calabi--Yau 簇 $X,Y$ 的导出范畴之间的任何等价都是 Fourier--Mukai 型的（从而 $\text{dim}(X)=\text{dim}(Y)$）；而命题 \ref{whenprojvarareequiv} 断言，若它们的凝聚层范畴等价，则 $X\cong Y$（但请注意其中第二个要点并不适用，因为 $K_X\cong\mathcal{O}_X$ 并不丰富！）；最后，$\text{dim}(X)=\text{dim}(Y)$ 蕴涵任何全忠实正合函子 $\mathsf{D^b}(X)\rightarrow\mathsf{D^b}(Y)$ 都是等价，此即命题 \ref{whenFourierMukaiisequivalence}。

\item 还有一些专属于 Calabi--Yau 簇的事实，可在 \cite{[Huy06]} 中找到：该书中猜想 6.24 预言两个\textit{双有理}的 Calabi--Yau 簇总是导出等价的（对 3-维簇已由推论 11.35 证实）；$\mathsf{D^b}(X)$ 不存在非平凡的\textit{半正交分解}（参习题 8.8）；$X$ 上任一线丛都是\textit{球状对象}（见 8.1 节），这一概念在镜像对称中屡屡出现；此外还有许多……
\end{itemize}

% [Tutor Note: 这份清单是全书对 $B$-侧的「验收表」，值得逐条勾稽。最容易被问倒的一处是「$K_X\cong\mathcal{O}_X$ 并不丰富」：命题 \ref{whenprojvarareequiv} 的第二要点要求典范（或反典范）线丛**丰富**，而 Calabi--Yau 的 $K_X$ 恰恰**平凡**，故该要点在最关心的一类流形上失效——这不是遗漏，而是 Bondal--Orlov 重建定理的真实边界。]

最后，我们已准备好以更深入的理解回望 Kontsevich 的猜想。我们以 \cite[猜想 7.26]{[Imp21]} 的形式陈述它，并略作改写。对 A-侧感兴趣之读者请参阅 \cite{[Imp21]}。

\begin{Con}[\textbf{同调镜像对称猜想}]\index{Homological Mirror Symmetry Conjecture}\label{HMSconj}		% HMS
	给定 Calabi--Yau 流形 $X^n=(X,J,\kappa^\mathbb{C},\varpi)$，存在同维数 $n$ 的「镜像」Calabi--Yau 流形 $\overline{X}=(\overline{X}, \overline{J},\overline{\kappa}^{\mathbb{C}},\overline{\varpi})$，以及三角范畴的等价
	\begin{equation}\label{HMS}
		\mathsf{D}^\pi(\widetilde{\mathscr{F}}(X))\cong \mathsf{D^b}(\overline{X}) \qquad\,\text{与}\qquad\, \mathsf{D}^\pi(\widetilde{\mathscr{F}}(\overline{X}))\cong \mathsf{D^b}(X)\,.
	\end{equation}
	因此，下述对应关系成立：
	\begin{itemize}[leftmargin=0.5cm]
		\renewcommand{\labelitemi}{\textendash}
		\item 由拉格朗日子流形搭成的分裂封闭扭曲复形 $\mathcal{L}^\bullet\in\textup{obj}\big(\mathsf{D}^\pi(\widetilde{\mathscr{F}}(X))\big)$，对应凝聚层的有界上链复形 $\overline{\mathcal{E}}^\bullet\in\textup{obj}(\mathsf{D^b}(\overline{X}))$；对称地，$\overline{\mathcal{L}}^\bullet\in\textup{obj}\big(\mathsf{D}^\pi(\widetilde{\mathscr{F}}(\overline{X}))\big)$ 对应 $\mathcal{E}^\bullet\in\textup{obj}\big(\mathsf{D^b}(X)\big)$。
		\item 对每个 $n\geq 0$，Floer 上同调 \textup{Hom}-群同构于凝聚 \textup{Ext}-群（作为 $\mathbb{C}$-向量空间），
		\begin{small}\begin{align*}
				& \textup{Hom}_{\mathsf{D}^\pi(\widetilde{\mathscr{F}}(X))}^n(\mathcal{L}_0^\bullet,\mathcal{L}_1^\bullet)\equiv HF^n(\mathcal{L}_0^\bullet,\mathcal{L}_1^\bullet)\cong\textup{Ext}_{\mathcal{O}_{\overline{X}}}^n(\overline{\mathcal{E}}_0^\bullet,\overline{\mathcal{E}}_1^\bullet)\equiv\textup{Hom}_{\mathsf{D^b}(\overline{X})}(\overline{\mathcal{E}}_0^\bullet,\overline{\mathcal{E}}_1^\bullet[n])\,, \\
				&\textup{Hom}_{\mathsf{D}^\pi(\widetilde{\mathscr{F}}(\overline{X}))}^n(\overline{\mathcal{L}}_0^\bullet,\overline{\mathcal{L}}_1^\bullet)\equiv HF^n(\overline{\mathcal{L}}_0^\bullet,\overline{\mathcal{L}}_1^\bullet)\cong\textup{Ext}_{\mathcal{O}_X}^n(\mathcal{E}_0^\bullet,\mathcal{E}_1^\bullet)\equiv\textup{Hom}_{\mathsf{D^b}(X)}(\mathcal{E}_0^\bullet,\mathcal{E}_1^\bullet[n])\,.
		\end{align*}\end{small}
		\!\!\!值得注意的是，$H^n(\mathcal{L}^\bullet;\Lambda_\mathbb{C})\!\cong\!\textup{Ext}_{\mathcal{O}_{\overline{X}}}^n(\overline{\mathcal{E}}^\bullet,\overline{\mathcal{E}}^\bullet)$ 且 $H^n(\overline{\mathcal{L}}^\bullet;\Lambda_\mathbb{C})\!\cong\!\textup{Ext}_{\mathcal{O}_{X}}^n(\mathcal{E}^\bullet,\mathcal{E}^\bullet)$。
		\item 在这些同构之下，Floer 积
		\begin{align*}
			& \langle \mu^2\rangle: HF(\mathcal{L}_1^\bullet,\mathcal{L}_2^\bullet)\otimes HF(\mathcal{L}_0^\bullet,\mathcal{L}_1^\bullet)\rightarrow HF(\mathcal{L}_0^\bullet,\mathcal{L}_2^\bullet)\quad\text{与} \\
			& \langle \mu^2\rangle: HF(\overline{\mathcal{L}}_1^\bullet,\overline{\mathcal{L}}_2^\bullet) \otimes HF(\overline{\mathcal{L}}_0^\bullet,\overline{\mathcal{L}}_1^\bullet)\rightarrow HF(\overline{\mathcal{L}}_0^\bullet,\overline{\mathcal{L}}_2^\bullet)
		\end{align*}
		分别化为 Yoneda 积
		\begin{align*}
			& \diamond:\textup{Ext}(\overline{\mathcal{E}}_1^\bullet,\overline{\mathcal{E}}_2^\bullet)\otimes\textup{Ext}(\overline{\mathcal{E}}_0^\bullet,\overline{\mathcal{E}}_1^\bullet) \rightarrow\textup{Ext}(\overline{\mathcal{E}}_0^\bullet,\overline{\mathcal{E}}_2^\bullet)\qquad\text{与} \\
			& \diamond:\textup{Ext}(\mathcal{E}_1^\bullet,\mathcal{E}_2^\bullet)\otimes\textup{Ext}(\mathcal{E}_0^\bullet,\mathcal{E}_1^\bullet) \rightarrow\textup{Ext}(\mathcal{E}_0^\bullet,\mathcal{E}_2^\bullet)\,.
		\end{align*}
	\end{itemize}
	从物理学的观点看，我们得到如下双射对应：$X$ 上、由 $(X,\kappa^\mathbb{C})$ 紧化的拓扑扭曲 \textup{A}-模型的稳定 \textup{A}-膜 $\mathcal{L}^\bullet$ 的模空间，与 $\overline{X}$ 上、由 $(\overline{X},\overline{J})$ 紧化的拓扑扭曲 \textup{B}-模型的稳定 \textup{B}-膜 $\overline{\mathcal{E}}^\bullet$ 的模空间之间的双射，反之亦然。
\end{Con}

\begin{Rem}
	\cite[注 7.27]{[Imp21]} 中的告诫依然适用。其中一个微妙的方面（那里已考虑到）是：Fukaya 范畴更一般地（也更正确地！）是在 Novikov 域 $\Lambda_\mathbb{K}$ 上定义的（见 \cite[定义 5.9]{[Imp21]} 或下文定义 \ref{Seidelstuff}），如此方可避免 Floer 复合映射中出现无限和。相应地，应当要求所猜想的镜像 Calabi--Yau 流形是 $\Lambda_\mathbb{K}$ 上的代数簇，而非基域 $\mathbb{K}$ 上的！所幸，在我们偏爱的复情形 $\mathbb{K}=\mathbb{C}$ 中，$\Lambda_\mathbb{C}$ 是代数闭的，因此其上的代数几何仍相当「标准」。在第 7 章处理椭圆曲线时，我们将（安全地）忽略这一问题，而 \cite{[AS09]} 对它作了彻底处理。\footnote{感谢 Biran 教授促使我着重指出这一非平凡的失洽之处。}
\end{Rem}

% [Tutor Note: 这条告诫是全文最微妙的技术性缺陷，作者本人也承认。要点：A-侧的 Fukaya 范畴必须定义在 Novikov 域 $\Lambda_{\mathbb{K}}=\{\sum a_i T^{\lambda_i}\}$ 上（$T$ 为形式参数、$\lambda_i$ 趋于 $+\infty$ 的实指数），以保证 Floer 复合的收敛；而 B-侧是普通域 $\mathbb{K}$ 上的代数几何。两边若要在同一基域上作范畴等价，就要求 $\Lambda_{\mathbb{K}}$ 与 $\mathbb{K}$ 相容——只在 $\mathbb{K}=\mathbb{C}$ 且取 $\Lambda_\mathbb{C}$ 时才免于此虑。椭圆曲线情形（第 7 章）之所以能绕开，是因为此时可用显式的 $A_\infty$-结构（Polishchuk--Zaslow）而不必诉诸一般理论。]

注意：由分裂封闭构造最终在 A-侧取上同调，猜想 \ref{HMSconj} 作用在经典范畴的层面上（而有界导出范畴直接就已是标准范畴，这一点我们非常清楚）。既然对 B-侧已有更牢固的把握，我们便为 \eqref{HMS} 提供一种额外的、隐含的解释，即把它看作 $A_\infty$-范畴之间的对偶（我们参阅 \cite[第 8.2.1 小节]{[ABC+09]}，并采用 \cite[第 2.2 节]{[Imp21]} 的记号）。

\begin{Rem}
	设 $X$ 是代数闭域 $\mathbb{K}$ 上的光滑代数簇。则我们可以按下述方式赋予其有界导出范畴一个 $A_\infty$-范畴的结构 $\mathsf{D^b_\infty}(X)$。令 $\text{obj}(\mathsf{D^b_\infty}(X))\coloneqq\text{obj}(\mathsf{D^b}(X))$，并对任意给定的 $\mathcal{F}^\bullet,\mathcal{G}^\bullet\in\text{obj}(\mathsf{D^b_\infty}(X))$ 定义 $\mathbb{Z}$-分次 $\mathbb{K}$-向量空间
	\begin{align*}
		\text{hom}_{\mathsf{D^b_\infty}\!(X)}(\mathcal{F}^\bullet,\mathcal{G}^\bullet)\coloneqq\; &\text{Hom}_{\mathcal{O}_X}^\bullet(\mathcal{F}^\bullet,\mathcal{G}^\bullet)\qquad\text{其中} \\ &\text{Hom}_{\mathcal{O}_X}^n(\mathcal{F}^\bullet,\mathcal{G}^\bullet)=\bigoplus_{k+l=n}\check{C}^k(\mathfrak{U},\mathcal{H}om_{\mathcal{O}_X}^l(\mathcal{F}^\bullet,\mathcal{G}^\bullet))\,,
	\end{align*}
	其中 Čech 复形如 \eqref{Cechcomplex} 所定义，$\mathcal{H}om$ 如 \eqref{locHomoncomplexes} 所定义，微分分别为 $\check{d}^\bullet$（如 \eqref{Cechdifferential}）与 $d^\bullet:\mathcal{H}om_{\mathcal{O}_X}^\bullet(\mathcal{F}^\bullet,\mathcal{G}^\bullet)\rightarrow \mathcal{H}om_{\mathcal{O}_X}^{\bullet+1}(\mathcal{F}^\bullet,\mathcal{G}^\bullet)$。
	\newline\big(后者更明确地由
	\[
	(d^n(\chi))_m=d_\mathcal{G}^{n+m}\circ \chi_m-(-1)^n\chi_{m+1}\circ d_\mathcal{F}^m\in\mathcal{H}om_{\mathcal{O}_X}(\mathcal{F}^m,\mathcal{G}^{n+m+1})
	\]
	给出，对任一 $\chi=(\chi_m)_{m\in\mathbb{Z}}\in\mathcal{H}om_{\mathcal{O}_X}^n(\mathcal{F}^\bullet,\mathcal{G}^\bullet)$，其中 $\chi_m\in\mathcal{H}om_{\mathcal{O}_X}(\mathcal{F}^m,\mathcal{G}^{n+m})$。\big)

	因此，任一（齐次的）$f\in\text{hom}_{\mathsf{D^b_\infty}\!(X)}^n(\mathcal{F}^\bullet,\mathcal{G}^\bullet)=\text{Hom}_{\mathcal{O}_X}^n(\mathcal{F}^\bullet,\mathcal{G}^\bullet)$（次数为 $n$）由一族 $(f_{n,i_0...i_k})$ 给定（其中 $1\leq i_0<...<i_k\leq |I|$，仍按 \eqref{Cechcomplex} 的约定），由此可知
	\begin{align*}
		&\mu_{\mathsf{D^b_\infty}\!(X)}^1:\text{hom}_{\mathsf{D^b_\infty}\!(X)}(\mathcal{F}^\bullet,\mathcal{G}^\bullet)\rightarrow\text{hom}_{\mathsf{D^b_\infty}\!(X)}(\mathcal{F}^\bullet,\mathcal{G}^\bullet)[1]\,, \\
		&\big(\mu_{\mathsf{D^b_\infty}\!(X)}^1(f)\big)_{n\!+\!1,i_0...i_{k+1}}\!\coloneqq\!(-1)^n(\check{d}^kf)_{n,i_0...i_{k+1}}\!+\!(-1)^{n+k} d^n(f)_{n\!+\!1,i_0...i_k}
	\end{align*}
	是 $\text{hom}_{\mathsf{D^b_\infty}\!(X)}(\mathcal{F}^\bullet,\mathcal{G}^\bullet)$ 上良定义的微分。

	此外，设 $g=(g_{m,j_0...j_{k'}})\in\text{hom}_{\mathsf{D^b_\infty}\!(X)}^{m}(\mathcal{G}^\bullet,\mathcal{E}^\bullet)$，我们可以定义二阶复合映射
	\begin{align*}
		&\mu_{\mathsf{D^b_\infty}\!(X)}^2:\text{hom}_{\mathsf{D^b_\infty}\!(X)}(\mathcal{G}^\bullet,\mathcal{E}^\bullet)\otimes\text{hom}_{\mathsf{D^b_\infty}\!(X)}(\mathcal{F}^\bullet,\mathcal{G}^\bullet)\rightarrow\text{hom}_{\mathsf{D^b_\infty}\!(X)}(\mathcal{F}^\bullet,\mathcal{E}^\bullet)\,, \\
		&\big(\mu_{\mathsf{D^b_\infty}\!(X)}^2(g,f)\big)_{n,i_0...i_{k+k'}}\!\!=\!\big((-1)^{n\!+\!k(m\!-\!k')}g_{2n-k,i_0...i_{k'}}\!\circ\! f_{n,i_{k'}...i_{k+k'}}\big)|_{U_{i_0}\cap...\cap U_{i_{k+k'}}}.
	\end{align*}

	最后，由于层的态射复合本身已是结合的，我们可以令 $\mu_{\mathsf{D^b_\infty}\!(X)}^d\coloneqq 0$（$d>2$）。可以证明 $\mu_{\mathsf{D^b_\infty}\!(X)}^1$ 与 $\mu_{\mathsf{D^b_\infty}\!(X)}^2$ 确实满足 $A_\infty$-结合性方程。于是我们成功地为 $\mathsf{D^b_\infty}(X)$ 构造了一个 $A_\infty$-结构。\footnote{由于更高阶的复合映射平凡，$\mathsf{D^b_\infty}(X)$ 不止是 $A_\infty$-范畴：它是 \textit{DG-范畴}（微分分次范畴）。}

	诸 $\text{hom}_{\mathsf{D^b_\infty}\!(X)}$-空间满足 $H^n(\text{hom}_{\mathsf{D^b_\infty}\!(X)}(\mathcal{F}^\bullet,\mathcal{G}^\bullet))\!=\!\text{Ext}_{\mathcal{O}_X}^n(\mathcal{F}^\bullet,\mathcal{G}^\bullet)$，此由 \eqref{ExtRHomforcomplexes}。事实上，可以定义包含 $\text{Ext}_{\mathcal{O}_X}^\bullet(\mathcal{F}^\bullet,\mathcal{G}^\bullet)\hookrightarrow \text{hom}_{\mathsf{D^b_\infty}\!(X)}(\mathcal{F}^\bullet,\mathcal{G}^\bullet)$，它们是拟同构，其逆为 $H^n$；这又使我们能给上同调范畴 $H(\mathsf{D^b_\infty}(X))$（见 \cite[定义 2.6]{[Imp21]}）赋予 $A_\infty$-范畴的结构。但关键事实是：作为普通范畴有 $H(\mathsf{D^b_\infty}(X))\cong\mathsf{D^b}(X)$，于是我们为标准的有限导出范畴获得了 $A_\infty$-结构（当心：其高阶复合映射非平凡！）。

	其寓意是：我们可以把猜想 \ref{HMSconj} 也理解为（$A_\infty$-范畴的零次上同调范畴之间的）一种对偶性！
\end{Rem}

% [Tutor Note: 这条注给出的是全书最深的一步：把 $B$-侧「升格」为 $A_\infty$-范畴，从而让猜测的两边处在同一类结构上，而不再只是「三角范畴等价」这种弱化陈述。几个易错点：(1) 微分 $\mu^1$ 中的两个符号因子 $(-1)^n$ 与 $(-1)^{n+k}$ 来自 Čech 微分与内部 Hom 微分的符号约定（即 Koszul 符号律），是极易写错的地方，改换 Čech 记号时须重推；(2) $\mu^2$ 公式中的 $g_{2n-k,\,i_0...i_{k'}}$ 下标 $2n-k$ 而非 $n+m$，这是把内部 Hom 的度数也计入后的结果；(3) 结论「$H(\mathsf{D^b_\infty}(X))\cong\mathsf{D^b}(X)$ 作为普通范畴」是 $A_\infty$-形式化的一般事实（同伦范畴即取零次上同调），但**高阶** $\mu^{\geq 3}$ 在 $H$ 上一般不为零，故 $H(\mathsf{D^b_\infty}(X))$ 与 $\mathsf{D^b}(X)$ 只在普通范畴层面相等。]

